\begin{remark*}
The diameter of a set is defined from all pairwise distances inside that set.
Mathematically, if \(S\) is a subset of a metric space \((X,d)\), then
\[
  \operatorname{diam}(S)
  =
  \sup \{ d(x,y) : x \in S,\ y \in S \}.
\]
As with point-to-set distance, it is useful in Lean to name the set being
supremized:
\[
\begin{gathered}
  \texttt{def diameterSet}\\
  \texttt{\ \ \ \ \{Point : Type u\}}\\
  \texttt{\ \ \ \ (M : MetricSpace Point)}\\
  \texttt{\ \ \ \ (S : Set Point) : Set Real :=}\\
  \texttt{\ \ \{ r : Real | exists x : Point, x in S and}\\
  \texttt{\ \ \ \ exists y : Point, y in S and r = M.distance x y \}.}
\end{gathered}
\]
Then the diameter itself is just
\[
\begin{gathered}
  \texttt{noncomputable def diameter}\\
  \texttt{\ \ \ \ \{Point : Type u\}}\\
  \texttt{\ \ \ \ (M : MetricSpace Point)}\\
  \texttt{\ \ \ \ (S : Set Point) : Real :=}\\
  \texttt{\ \ sSup (diameterSet M S).}
\end{gathered}
\]
\end{remark*}

\begin{remark*}
The monotonicity theorem for diameter is an instance of a more general theorem:
supremum is monotone under set inclusion. The mathematical proof has two
layers.

First, set inclusion of the original sets gives set inclusion of the pairwise
distance sets. If \(A \subseteq B\), then any distance \(d(x,y)\) with
\(x,y \in A\) is also a distance with \(x,y \in B\). In Lean this is:
\[
\begin{gathered}
  \texttt{theorem diameterSet\_mono}\\
  \texttt{\ \ \ \ \{Point : Type u\}}\\
  \texttt{\ \ \ \ (M : MetricSpace Point)}\\
  \texttt{\ \ \ \ \{A B : Set Point\}}\\
  \texttt{\ \ \ \ (set\_inclusion : A subset B) :}\\
  \texttt{\ \ \ \ diameterSet M A subset diameterSet M B := by}\\
  \texttt{\ \ intro r r\_in\_diameter\_A}\\
  \texttt{\ \ rcases r\_in\_diameter\_A with}\\
  \texttt{\ \ \ \ <x, x\_in\_A, y, y\_in\_A, rfl>}\\
  \texttt{\ \ exact <x, set\_inclusion x\_in\_A,}\\
  \texttt{\ \ \ \ y, set\_inclusion y\_in\_A, rfl>.}
\end{gathered}
\]
The proof reads exactly like the hand proof. Start with an arbitrary member
\texttt{r} of \texttt{diameterSet M A}. By definition, it is
\texttt{M.distance x y} for some \texttt{x} and \texttt{y} in \texttt{A}. The
hypothesis \texttt{set\_inclusion} converts \texttt{x in A} into
\texttt{x in B}, and converts \texttt{y in A} into \texttt{y in B}. Therefore
the same real number \texttt{r} belongs to \texttt{diameterSet M B}.

Second, apply the general supremum theorem:
\[
  \texttt{csSup\_le\_csSup}.
\]
For real numbers, this theorem says: if \(s \subseteq t\), \(s\) is nonempty,
and \(t\) is bounded above, then
\[
  \sup s \le \sup t.
\]
The diameter theorem is then:
\[
\begin{gathered}
  \texttt{theorem diameter\_monotone\_under\_inclusion}\\
  \texttt{\ \ \ \ \{Point : Type u\}}\\
  \texttt{\ \ \ \ (M : MetricSpace Point)}\\
  \texttt{\ \ \ \ \{A B : Set Point\}}\\
  \texttt{\ \ \ \ (set\_inclusion : A subset B)}\\
  \texttt{\ \ \ \ (A\_diameterSet\_nonempty : (diameterSet M A).Nonempty)}\\
  \texttt{\ \ \ \ (B\_diameterSet\_bddAbove : BddAbove (diameterSet M B)) :}\\
  \texttt{\ \ \ \ diameter M A <= diameter M B := by}\\
  \texttt{\ \ exact csSup\_le\_csSup}\\
  \texttt{\ \ \ \ B\_diameterSet\_bddAbove}\\
  \texttt{\ \ \ \ A\_diameterSet\_nonempty}\\
  \texttt{\ \ \ \ (diameterSet\_mono M set\_inclusion).}
\end{gathered}
\]
The helper theorem \texttt{diameterSet\_mono} does the metric-space-specific
work: it turns \(A \subseteq B\) into inclusion of the two distance sets. The
general theorem \texttt{csSup\_le\_csSup} then does the order-theoretic work:
it turns inclusion of real sets into an inequality between their suprema.

This separation is useful. If another construction is also defined as the
supremum of a set of real values, it can reuse \texttt{csSup\_le\_csSup}; only
the construction-specific set-inclusion lemma has to be proved.
\end{remark*}
